% =============================================================================
% Appendix A. Moved out of 4.1 by Ben, 2026-10-18 ("let's put in appendix 4.1.4
% for now, I don't convinced it's valuable"). Single-node runs only, P/D
% emulated, so this block has not been compared against P/D orchestration.
% =============================================================================

\section{Decoupling Execution from Residency}\label{sec:residency}

% Figure -- Block A: decoupling execution from residency. Column width
% (3.35in = 8.5cm), drawn at final size, never \resizebox'ed.
%
% What the code does (perblock-d2h-opt):
%   (1) scheduler.py `_rungd_admit_preloaded` -> admitted with
%       `empty_kv_cache_blocks`; hybrid_rungd_pool.py `ensure_req` gives it one
%       contiguous extent of host block ids, freed only by `free_req`.
%   (2) scheduler.py `_rungd_defer_for_riders`: a block-starved request is
%       stepped over (it keeps its place) when a preloaded rider waits behind.
%   (3) scheduler.py `_rungd_reclaim_seeded` -> kv_cache_manager
%       `reclaim_tier_blocks`, only on the worker's seed acknowledgement.
\begin{figure}[t]
\centering
\begin{tikzpicture}[x=1cm,y=1cm]
\tikzset{num/.style={circle, fill=black, text=white, inner sep=0.4pt,
         minimum size=3.6mm, font=\sffamily\bfseries\scriptsize},
         blk/.style={cell, minimum width=4.6mm, minimum height=4.6mm,
                     font=\sffamily\tiny},
         req/.style={draw, rounded corners=1.5pt, line width=0.6pt,
                     minimum height=0.50cm, font=\sffamily\scriptsize, align=center}}

%% ----------------------------------------------------- HBM: paged block space
\fill[hbmF!45, rounded corners=3pt] (0,2.90) rectangle (8.45,4.55);
\node[hd, text=hbmD, anchor=north west] at (0.08,4.52) {HBM {\mdseries\scriptsize paged block space}};
\foreach \i/\o in {0/R1,1/R1,2/R2,3/R2,4/R2}
  \node[blk, fill=hbmF, draw=hbmD] at (0.40+\i*0.50,3.85) {\o};
\foreach \i in {5,...,12}
  \node[blk, fill=white, draw=black!30] at (0.40+\i*0.50,3.85) {};
\node[sub, anchor=south] at (4.40,4.10) {free list};
% blocks that were R3's: reclaimed after the ack, reissued to R4
\foreach \i in {13,14,15}
  \node[blk, fill=hbmF, draw=hbmD, line width=0.8pt] (h\i) at (0.40+\i*0.50,3.85) {R4};
\draw[decorate, decoration={brace, amplitude=1.2mm, mirror}, hbmD]
  (6.80,3.55) -- (8.00,3.55);
\node[font=\sffamily\scriptsize, anchor=north east, align=right] at (7.05,3.42)
  {were R3's: \textbf{reclaimed} after the ack, reissued to R4};
\node[num] at (1.20,3.27) {3};

%% ----------------------------------------------------- scheduler: waiting queue
\fill[ctlF!60, rounded corners=3pt] (0,1.60) rectangle (8.45,2.75);
\node[hd, text=ctlD, anchor=north west] at (0.08,2.72) {Waiting queue};
\node[req, draw=hbmD, fill=hbmF, minimum width=1.65cm] (r4) at (6.90,2.02)
  {R4 {\tiny(head)}\\[-2pt]{\tiny needs blocks}};
\node[req, draw=ddrD, fill=ddrF, minimum width=1.65cm] (qa) at (4.55,2.02)
  {rider A\\[-2pt]{\tiny needs 0 blocks}};
\node[sub] at (2.70,2.02) {$\cdots$};
\draw[flow, ctlD, line width=0.8pt] ($(r4.north west)+(0.25,0)$)
  .. controls +(-0.3,0.30) and +(0.4,0.30) .. ($(qa.north east)+(-0.25,0)$);
\node[sub, anchor=south] at (5.72,2.50) {step over};
\node[num] at (3.40,2.02) {2};
% R4 is then served from the reissued blocks
\draw[flow, hbmD, line width=0.8pt] (r4.north -| h14) -- (h14.south);

%% ----------------------------------------------------- host DDR: extent space
\fill[ddrF!50, rounded corners=3pt] (0,-0.40) rectangle (8.45,1.40);
\node[hd, text=ddrD, anchor=north west] at (0.08,1.37) {Host DDR {\mdseries\scriptsize private extent space}};
% extent = one contiguous run of host block ids, sized for the request
\newcommand{\extent}[5]{% x0, used, total, label, name
  \foreach \j in {1,...,#2}
    \node[blk, fill=ddrF, draw=ddrD] at (#1+\j*0.46-0.23,0.22) {};
  \pgfmathtruncatemacro{\nx}{#2+1}
  \foreach \j in {\nx,...,#3}
    \node[blk, fill=white, draw=ddrD!50] at (#1+\j*0.46-0.23,0.22) {};
  \draw[ddrD, line width=0.5pt, rounded corners=1pt]
    (#1-0.03,-0.03) rectangle (#1+#3*0.46+0.03,0.47);
  \node[font=\sffamily\scriptsize, anchor=north] at (#1+#3*0.23,-0.05) {#4};
  \coordinate (#5) at (#1+#3*0.23,0.47);
}
\extent{0.20}{3}{5}{request A}{eA}
\extent{2.70}{2}{5}{request B}{eB}
\extent{5.20}{4}{6}{request R3 (demoted)}{eR3}

%% ----------------------------------------------------- the movements
% (1) rider A admitted with zero HBM blocks; its KV is in its own extent
\draw[flow, ddrD, line width=0.8pt] (qa.south) -- (qa.south |- 0,0.85) -| (eA);
\node[num] at (2.40,0.85) {1};
% (3) R3 demoted: KV seeded into extent 2; the ack precedes the reclaim
\draw[d2h] (h15.east) -- (8.28,3.85) -- (8.28,0.85) -| (eR3);
\node[font=\sffamily\scriptsize, text=d2hC, anchor=east] at (8.20,1.05) {seed, then ack};

\end{tikzpicture}
\caption{\textbf{A rider runs without holding HBM.}
\stepn{1}~A rider is admitted with an \emph{empty} device block list: its KV
already sits in a private, contiguous host extent, which it keeps until it
finishes. \stepn{2}~A request starved of blocks no longer blocks the waiting
queue for riders behind it, which need none. \stepn{3}~Demoting resident R3
copies its KV into an extent; only on the worker's acknowledgement are its HBM
blocks reclaimed and reissued (here, to R4). Host extents are keyed by request,
not by device block id, so reissue cannot reach a rider's copy. Measured at a fixed
KV pool: 48 concurrent requests against an all-HBM ceiling of 8 (synthetic
rider KV, capacity only); 133 of 192 blocks reclaimed on demotion.
\torun{single node with emulated P/D handover; no comparison against other P/D
orchestration has been run.}}
\label{fig:residency}
\end{figure}


\paragraph{Challenge.} In a paged engine, being scheduled means holding device
blocks: the waiting queue admits a request by allocating its blocks, and the
running batch is bounded by them. A rider must be schedulable while holding none.

\paragraph{Design.} Three scheduler changes and one address space
(\cref{fig:residency}). A rider is admitted with an empty block list, and only
once its KV, written by the prefill node into host DDR, is registered in a host
extent; the worker refuses a rider whose KV was never delivered. The waiting
queue no longer stops at the first request it cannot allocate, so a
block-starved request at its head does not stall a rider behind it. A resident
is demoted through an explicit acknowledgement, because the scheduler cannot
observe a host-side write: its KV is copied into an extent, the worker
acknowledges it, and only then are its device blocks reclaimed. And the host
copy is addressed by a private, per-request extent space rather than by device
block id: a network-admitted rider has no device id to be keyed by, and a
demoted rider's ids are returned to the free list and reissued to other
requests.

\paragraph{Capacity.} At a fixed KV pool whose all-HBM ceiling is 8 concurrent
requests, the node held 48 with riders on top of a full resident cohort, and
demotion reclaimed 133 of 192 device blocks. The capacity run fills rider KV
synthetically, so it measures capacity only; token correctness of the same path
is tested bitwise and separately.
% src: benchmarks/kv_pipelining_bench/hybrid_preload_capacity.py docstring and
%      ARCHITECTURE.md s9 (8 -> 48 at b_hbm_max = measured ceiling);
%      hybrid_reclaim_correctness.py (133 of 192)
\torun{all runs are single-node with the P/D handover emulated; this block has
not been compared against other P/D orchestration.}
